% III. МЕТОДИКА И ЕКСПЕРИМЕНТАЛНИ РЕЗУЛТАТИ, част 1: ПЛАНЪТ.
% ВНИМАНИЕ: тук няма и не трябва да има нито едно измерено число. Числата се
% появяват едва в chapters/06_results.tex, което е продължение на СЪЩИЯ раздел.
\section{Методика и експериментални резултати}
\label{sec:method}

\subsection{Методика}

Проверката отговаря на четири отделни въпроса: дали извлеченият модел отговаря на изходния текст, колко
струва извличането, дали пълният цикъл запазва модела, и дали проверяващият модул открива нарочно
внесено разминаване. Смесването им в едно число би скрило точно това, което трябва да бъде видяно.
Измерванията са на Windows 10.0.26200, AMD Ryzen 5 3600 с 6 ядра, построяване с GNU 15.2.0
(MinGW-w64) в режим Release, CMake 4.3.2 с Ninja 1.13.2, стандарт C++20, а фронтендът при извличането
е clang 22.1.0. Средата се извежда от измервателния скрипт, тъй като резултатите зависят от версията
на фронтенда. Всяко измерване следва чисто построяване, времената се повтарят по три пъти и се
съобщава \term{най-малкото}, а не средното: по-бавните изпълнения се различават по планиране на
процесите, което не е свойство на инструмента. Всичко се възпроизвежда с \code{cmake --build build
--target measure}, а четвъртият експеримент с целта \code{divergence}.

При \textbf{първия експеримент} очакваният модел за всяка конструкция е записан на ръка \term{преди}
инструментът да бъде пуснат върху нея, тоест сравнението не е между два изхода на самия инструмент;
отделно върху породените заготовки се пуска фронтендът с \code{-fsyntax-only}, защото пораждащата
посока пише код, който никой човек не е чел. \textbf{Вторият} използва породени единици с нарастващ
брой класове, всеки с по един представител на всяка форма на член от Табл.~\ref{tab:ownership}, и двата
действителни входа; фронтендът, разборът и обхождането се измерват поотделно, защото се променят
различно. \textbf{Третият} превръща текст в модел, модела в диаграма и диаграмата обратно в модел;
показателят е дял на \term{запазените елементи}, а не единица минус дял на разликите, защото когато цял
класификатор не се върне, членовете му изобщо не се сравняват. \textbf{Четвъртият} променя само едната
страна на съгласувана двойка, по една промяна наведнъж, като преди всяка промяна се проверява, че
недокоснатата двойка съвпада; показателите са два, откриваемост и брой съобщения извън очакваното,
защото инструмент, който съобщава за всичко, постига идеална откриваемост и е безполезен.
